
\examquestion{2023}{2}{Hydrogen radial probabilities, helium spins and E1 lines}{3}
The hydrogen \(1s\) wavefunction is
\[
\psi_{1s}(r)=\frac{e^{-r/a_0}}{\sqrt{\pi a_0^3}},
\qquad
|\psi_{1s}|^2=\frac{e^{-2r/a_0}}{\pi a_0^3}.
\]

\subsection*{(a) Radius with maximum detection probability}
The correct \emph{radial probability distribution}
is the probability in a spherical shell,
\(P(r)\diff r=4\pi r^2|\psi|^2\diff r\):
\[
P(r)=\frac{4r^2}{a_0^3}e^{-2r/a_0}.
\]
Differentiating the log,
\(P'/P=2/r-2/a_0\). Apart from the zero
endpoint, the maximum is
\[
\boxed{r_{\rm most\ probable}=a_0}.
\]
Maximising \(|\psi(r)|^2\) alone would wrongly give
\(r=0\): it omits the increasing spherical-shell volume.

\subsection*{(b) Evaluate \(\langle r\rangle\) and
\(\langle1/r\rangle\)}
Using
\(\int_0^\infty r^n e^{-\beta r}\diff r
=n!/\beta^{n+1}\),
\begin{align*}
\langle r\rangle
&=\frac4{a_0^3}\int_0^\infty
r^3e^{-2r/a_0}\diff r
=\frac4{a_0^3}\frac{3!}{(2/a_0)^4}
=\boxed{\frac{3a_0}{2}},\\
\left\langle\frac1r\right\rangle
&=\frac4{a_0^3}\int_0^\infty
re^{-2r/a_0}\diff r
=\frac4{a_0^3}\frac{1!}{(2/a_0)^2}
=\boxed{\frac1{a_0}}.
\end{align*}
Bohr's \(n=1\) orbit has fixed radius \(a_0\);
it agrees with the quantum \emph{most probable}
radius and \(\langle1/r\rangle=1/a_0\),
but not the mean radius \(\langle r\rangle=3a_0/2\).

\subsection*{(c) Schr\"odinger equation for helium}
For a fixed nucleus \(Z=2\), a simplified
non-relativistic two-electron Hamiltonian is
\[
\boxed{H=
-\frac{\hbar^2}{2m_e}(\nabla_1^2+\nabla_2^2)
-\frac{2e^2}{4\pi\epsilon_0}
\left(\frac1{r_1}+\frac1{r_2}\right)
+\frac{e^2}{4\pi\epsilon_0r_{12}}.}
\]
The three terms represent electron kinetic
energy, attractions to the nucleus and mutual
electron repulsion, respectively. Solve
\(H\Psi(\mathbf r_1,\mathbf r_2)=E\Psi\)
with total fermion-exchange antisymmetry.
Spin and relativistic interactions would be
additional smaller corrections.

\subsection*{(d) All four two-electron spin states}
\[
\begin{array}{c|l|c}
\text{Total }S&\text{Spin state}&\text{Exchange}\\\hline
1,\ M_S=1&\ket{\uparrow\uparrow}&\text{symmetric}\\
1,\ M_S=0&
(\ket{\uparrow\downarrow}+\ket{\downarrow\uparrow})/\sqrt2&
\text{symmetric}\\
1,\ M_S=-1&\ket{\downarrow\downarrow}&\text{symmetric}\\
0,\ M_S=0&
(\ket{\uparrow\downarrow}-\ket{\downarrow\uparrow})/\sqrt2&
\text{antisymmetric}
\end{array}
\]
The first three are the triplet, the last the singlet.
Since total two-electron wavefunction is antisymmetric,
a triplet needs an antisymmetric spatial part, while
a singlet needs a symmetric spatial part.

\subsection*{(e) Physical origin and relative magnitude of hyperfine structure}
Fine structure primarily concerns electronic
spin--orbit and relativistic effects. Hyperfine
structure couples electronic angular momentum and
fields to \emph{nuclear} magnetic dipole/electric
quadrupole moments. The nuclear magnetic moment
is typically of nuclear-magneton scale,
\(\mu_N=e\hbar/(2m_p)\), much smaller than the
electron Bohr magneton \(\mu_B=e\hbar/(2m_e)\);
this suppresses hyperfine interactions.
For \({}^4\mathrm{He}\), whose nucleus has \(I=0\),
magnetic-dipole hyperfine splitting is absent.

\subsection*{(f) Allowed E1 transitions among the stated terms}
To leading order in LS coupling, E1 requires
opposite orbital parity, \(\Delta S=0\), and
\(|\Delta J|\le1\) with \(J=0\not\to J'=0\).
Thus the allowed combinations are
\[
\boxed{{}^1P_1\longleftrightarrow{}^1S_0,\qquad
{}^3P_{0,1,2}\longleftrightarrow{}^3S_1.}
\]
The direction of spontaneous emission is from the
higher to the lower member of each pair.
Singlet--triplet lines are spin-forbidden in pure
LS coupling but can gain weak intensity through
spin--orbit mixing. A simple term diagram is:
\[
\begin{tikzpicture}[x=1cm,y=.5cm]
\draw[-{Latex[length=2mm]}](0,0)--(0,6.7) node[above]{energy};
\draw[accent,thick](1,5)--(3.3,5) node[right,ink]{\({}^1P_1\)};
\draw[accent,thick](1,1.5)--(3.3,1.5) node[right,ink]{\({}^1S_0\)};
\draw[-{Latex[length=2mm]}](2,4.8)--(2,1.7);
\draw[accent,thick](6,5.5)--(8,5.5) node[right,ink]{\({}^3P_2\)};
\draw[accent,thick](6,4.9)--(8,4.9) node[right,ink]{\({}^3P_1\)};
\draw[accent,thick](6,4.3)--(8,4.3) node[right,ink]{\({}^3P_0\)};
\draw[accent,thick](6,1.4)--(8,1.4) node[right,ink]{\({}^3S_1\)};
\foreach \y/\x in {5.5/6.5,4.9/7,4.3/7.5}
{\draw[-{Latex[length=1.7mm]}](\x,\y-.2)--(7,1.6);}
\end{tikzpicture}
\]
\solcheck{Radial expectations require the
\(4\pi r^2\) measure. For selection rules,
check \emph{spin}, \emph{parity} and \emph{angular
momentum} independently.}
